% vlf_reception_notes.tex
%
% A short supplement for the docs/ folder of a SuperSID station.
% Build with:  pdflatex vlf_reception_notes.tex
% Twice, if you want the references resolved.

\documentclass[11pt,a4paper]{article}

\usepackage[T1]{fontenc}
\usepackage[utf8]{inputenc}
\usepackage{amsmath}
\usepackage{amssymb}
\usepackage[margin=2.5cm]{geometry}
\usepackage{natbib}
\usepackage{enumitem}
\usepackage{hyperref}
\hypersetup{colorlinks=true, linkcolor=black, citecolor=black, urlcolor=blue}

\setlength{\parskip}{0.6em}
\setlength{\parindent}{0pt}

\title{VLF reception, station geometry and hourly filtering}
\author{Notes for a SuperSID monitoring station}
\date{\today}

\begin{document}
\maketitle

\section{Induced voltage in the loop antenna}

A vertically polarized wave in the waveguide carries a horizontal magnetic
field that is perpendicular to its direction of travel \citep{loudet2013}. The
voltage induced in the loop is:

\begin{equation}
     \left [\frac{V_{\mathrm{rms}}}{\mathrm{V}} \right ]= 2\pi \left[ \frac{f}{\mathrm{Hz}} \right] N \left[ \frac{A_{\mathrm{eff}}}{\mathrm{m}^2} \right] \left[ \frac{B}{\mathrm{T}} \right] \cos\theta
    \label{eq:faraday}
\end{equation}

where:
\begin{description}[itemsep=2pt,topsep=2pt]
    \item[$-$ $V_{\mathrm{rms}}$] is the root mean square voltage across the loop.
    \item[$-$ $f$] is the transmitter frequency.
    \item[$-$ $N$] is the number of turns in the loop.
    \item[$-$ $A_{\mathrm{eff}}$] is the effective area enclosed by one turn.
    \item[$-$ $B$] is the magnetic flux density of the arriving wave.
    \item[$-$ $\theta$] is the angle between the loop axis and the direction to
        the transmitter, so the response falls to zero when the loop faces it
        edge on.
\end{description}

The $\cos\theta$ term is why the bearing to each transmitter matters when the
antenna is put up, and the two sections below are how that bearing and the path
length are worked out.

\section{Forward azimuth to a transmitter}

\begin{equation}
    \left [\frac{\theta}{\mathrm{rad}}\right] = \arctan \left( \frac{\sin(\Delta\lambda)}{\cos(\phi_1)\tan(\phi_2) - \sin(\phi_1)\cos(\Delta\lambda)} \right)
    \label{eq:azimuth}
\end{equation}

where:
\begin{description}[itemsep=2pt,topsep=2pt]
    \item[$-$ $\theta$] is the forward azimuth in radians.
    \item[$-$ $\phi_1$] is the latitude of the starting point.
    \item[$-$ $\phi_2$] is the latitude of the destination point.
    \item[$-$ $\Delta\lambda$] is the difference in longitude between the two points.
\end{description}

\section{Great circle distance}

{\small
\begin{equation}
    \left [\frac{d}{\mathrm{km}} \right]= \left[ \frac{R_{E}}{\mathrm{km}} \right] \arccos\big(\sin\phi_{1}\sin\phi_{2} + \cos\phi_{1}\cos\phi_{2}\cos(\lambda_{2}-\lambda_{1})\big)
    \label{eq:gc}
\end{equation}
}

where:
\begin{description}[itemsep=2pt,topsep=2pt]
    \item[$-$ $d$] is the Great Circle distance.
    \item[$-$ $R_{E}$] is the mean radius of the Earth ($\approx 6371 \, \mathrm{km}$).
    \item[$-$ $\phi_{1}, \phi_{2}$] are the latitudes of the two points.
    \item[$-$ $\lambda_{1}, \lambda_{2}$] are the longitudes of the two points.
\end{description}

Equations~\ref{eq:azimuth} and~\ref{eq:gc} are what \texttt{vlf\_transit.py}
evaluates for every transmitter in its table, which is how the installer can
offer the ten or twenty nearest ones and how \texttt{vlf\_table.pdf} lists a
bearing and a distance for each.

\section{Why the hourly figure says BEMA was applied here}

The plot label reads \texttt{BEMA wing 6 (applied here, the csv itself is raw)}
because the smoothing seen in the figure is performed by the plotting script,
not by SuperSID. Although \texttt{log\_type} may be set to \texttt{filtered} in
\texttt{supersid.cfg}, that setting governs only the files written at the end of
each UTC day. The hourly dump produced by \texttt{hourly\_save} is written with
\texttt{log\_type='raw'} hardcoded in \texttt{supersid.py}, so the values in
\texttt{hourly\_current\_buffers.*.csv} are exactly as captured from the sound
card, and the file's own header confirms this with the line
\texttt{\# LogType = raw}.

To make the hourly figure comparable with the finished daily record, the script
reproduces SuperSID's BEMA filter (\texttt{SidFile.filter\_buffer}) on the data
before drawing: for each sample it takes the lowest value within a window of
$2 \times \mathit{bema\_wing}$ points, then applies a single moving average of
$2 \times \mathit{bema\_wing} + 1$ points over those minima. With the default
wing of 6 this spans roughly one minute at a five second log interval. Taking
the minimum first causes the curve to follow the quiet background level, so a
sudden ionospheric disturbance stands out as a departure from that baseline
rather than being averaged away. Where a file already carries
\texttt{LogType = filtered}, the script detects this from the header and applies
nothing further, and the label changes accordingly.

\subsection*{Shorter version, for use as a figure caption}

Signal from the hourly buffer, which SuperSID writes unfiltered regardless of
the \texttt{log\_type} setting. BEMA (wing 6: rolling minimum over 12 samples,
then a 13 sample moving average) is applied by the plotting script to match the
filtering used in the daily files.

\begin{thebibliography}{9}
\bibitem[Loudet(2013)]{loudet2013}
Loudet, L. (2013).
\newblock \emph{SID monitoring station}.
\newblock \url{https://sidstation.loudet.org/home-en.xhtml}
\end{thebibliography}

\end{document}
